% Study sheet: Volume II, Chapter 5, Collective Communication.
% Numbers: \val{...} from collective_communication_studysheet_numbers.tex, written by collective_communication_studysheet_numbers.py.
\documentclass{studysheet}
\input{collective_communication_studysheet_numbers}
\sheetvolume{2}
\sheetchapter{5}{Collective Communication}
\sheetsources{Numbers from mlsysim and the chapter's scenarios; KB, MB, GB decimal; assumed inputs marked.}

\begin{document}
% ===================================================================== FRONT: LEARN
\frontside
\begin{bigidea}
Every split of a model puts data on the wire. A \hl{collective} names what each rank holds; its
\hl{schedule} sets the startup rounds and bytes paid. In C\textsuperscript{3} terms, this prices the
Fleet Law's communication term.
\end{bigidea}

\begin{leftcol}
\sheetsec{Key ideas}
\begin{ideas}
\idea{A collective is a contract on results.}{AllReduce is ReduceScatter then AllGather, each half its
  bytes. Call the narrowest one the next step consumes.}
\idea{Price every round on its chunk.}{Each step pays startup plus chunk over bandwidth: big
  buckets on big groups can be latency bound.}
\idea{Rings save bytes, trees save rounds.}{The simple tree uses few rounds but resends the whole
  tensor each level.}
\idea{Hierarchy buys parallel ports.}{Reduce in the node, cross nodes with one shard per GPU,
  gather. On one shared uplink, the $\AcceleratorsPerNode$ shard rings run in turn.}
\idea{AllToAll ends at its active bound.}{Startup, injection, reception or the cut; under skew,
  the hottest receiver sets the layer.}
\idea{Shrink or hide, then judge by time to quality.}{Compression cuts only bytes; overlap hides network
  time, not processor overhead.}
\end{ideas}

\vspace{10pt}
\sheetsec{Which primitive moves it?}
\begin{sheettable}{@{}L l@{}}
\thead{Split (what moves)} & \thead{Primitive}\\ \midrule
Data parallel (full gradient) & AllReduce\\
Sharded, FSDP (shards) & AllGather, ReduceScatter\\
Tensor parallel (partial activations) & AllReduce\\
Pipeline (activations) & Send/Recv\\
Experts, embeddings (tokens) & AllToAll\\
\end{sheettable}

\vspace{10pt}
\fcard{Compression pays when}%
  {\CodecOverheadTime < \CommunicationTime(\DeviceCount)\left(1-\frac{1}{\CompressionRatio}\right),
   \qquad \ErrorFeedbackResidual_{\OptimizerStepIndex+1} = \LocalGradient_{\OptimizerStepIndex}
   + \ErrorFeedbackResidual_{\OptimizerStepIndex} - \CompressedGradientUpdate_{\OptimizerStepIndex}}%
  {Error feedback sends $\CompressedGradientUpdate$
   and keeps the residual $\ErrorFeedbackResidual$ for the next gradient $\LocalGradient$.}{}{}
\end{leftcol}%
\begin{rightcol}
\sheetsec{Formulas you will use}
\fcard{Ring AllReduce, priced per round}%
  {\RingAllReduceCompletionTime = \under{2(\DeviceCount-1)\MessageStartupLatency}{rounds}
   + \under{2\frac{\DeviceCount-1}{\DeviceCount}\frac{\CollectivePayloadSize}{\LinkBandwidth}}{bytes}}%
  {Each step sends a chunk $\CollectivePayloadSize/\DeviceCount$; below
   $\LatencyBandwidthCrossoverSize = \MessageStartupLatency\LinkBandwidth$ it is latency bound. Fusing
   $\SplitMessageCount$ small messages pays $\MessageStartupLatency$ once, not $\SplitMessageCount$ times.}%
  {NDR port: $\MessageStartupLatency = \val{alpha}\,\MicroUnitPrefix$s, $\LinkBandwidth = \val{beta}$~GB/s
   (\val{gbit}~Gb/s), so $\LatencyBandwidthCrossoverSize = \val{nstar}$~KB. Fit both on a benchmark
   matching operation, size, ranks, algorithm and placement (GPU, NIC, rail).}{}
\fcard{Simple tree and the crossover}%
  {\TreeAllReduceCompletionTime \approx 2\log_2\DeviceCount\left(\MessageStartupLatency
   + \frac{\CollectivePayloadSize}{\LinkBandwidth}\right), \qquad
   \RingTreeCrossoverPayload \approx \frac{\DeviceCount\MessageStartupLatency\LinkBandwidth}{\log_2\DeviceCount}}%
  {Roughly, ring wins above $\RingTreeCrossoverPayload$, the tree below. Double binary tree: tree rounds,
   ring bytes, $2\log_2\DeviceCount\,\MessageStartupLatency + 2\CollectivePayloadSize/\LinkBandwidth$.}{}{}
\fcard{Hierarchical AllReduce}%
  {\text{cross-node} = 2(\ClusterNodeCount-1)\left(\MessageStartupLatency
   + \frac{\CollectivePayloadSize}{\AcceleratorsPerNode\,\ClusterNodeCount\,\LinkBandwidth}\right)}%
  {Per GPU, a ring over the $\ClusterNodeCount$ nodes on its own port.
   Each NVLink phase moves $(\AcceleratorsPerNode-1)/\AcceleratorsPerNode$ of
   $\CollectivePayloadSize$ per GPU. A flat ring runs at one port: $2\CollectivePayloadSize/\LinkBandwidth$.}{}%
  {\begin{tikzpicture}[sheet]
     \pgfmathsetmacro\sc{200/\val{drawflat}}
     \node[note, anchor=west] at (0,40) {flat ring};
     \draw[nbox] (44,34) rectangle ({44+\val{drawflat}*\sc},46);
     \node[lab] at ({44+\val{drawflat}*\sc/2},40) {\val{flat} ms, one NDR port};
     \pgfmathsetmacro\a{44+\val{drawrs}*\sc} \pgfmathsetmacro\b{\a+\val{drawinter}*\sc}
     \pgfmathsetmacro\c{\b+\val{drawag}*\sc}
     \node[note, anchor=west] at (0,16) {hierarchical};
     \draw[accbox] (44,10) rectangle (\a,22);
     \draw[hotbox] (\a,10) rectangle (\b,22);
     \draw[accbox] (\b,10) rectangle (\c,22);
     \node[lab, anchor=west] at ({\c+4},16) {\val{hiertotal} ms, \val{hiergain}$\times$ faster};
     \node[small, anchor=north west] at (44,6) {NVLink \val{hierrs} + all ports \val{hierinter} + NVLink \val{hierag} ms};
     \node[note, anchor=south west] at (44,50) {\val{hiergb} GB on \val{hiernodes} nodes of \val{hiergpus} GPUs};
   \end{tikzpicture}}
\fcard{Overlap leaves a tail}%
  {\OverlappedBackwardCommunicationTime = \BackwardPassTime + \LayerAllReduceTime
   + (\LayerCount-1)\,\LayerExposedCommunicationTime}%
  {Over $\LayerCount$ layers, the last reduction is exposed and each excess
   $\LayerExposedCommunicationTime = \max(0,\,\LayerAllReduceTime-\LayerBackwardTime)$ piles up.}{}{}
\end{rightcol}

% ===================================================================== BACK: PRACTICE
\backside
\begin{answers}
\ans{Finish}{\val{treerounds} $\times$ \val{perround} = \val{treebytes} ms of bytes, \val{treetotal} ms in all:
  \val{treevsring}$\times$ slower than the ring's \val{ringlarge} ms, as \val{bucket} MB is above the \val{crossover} MB crossover. Tree rounds,
  ring bytes: \val{treestartup} + \val{byteslarge} = \val{doubletree} ms, \val{doubletreegain}$\times$ faster
  than the ring (tree card).}
\ans{1}{The full sum, mostly discarded. ReduceScatter leaves each rank its reduced shard for half the
  bytes (key idea one).}
\ans{2}{(a) Twice \val{qgb} GB $\div$ \val{beta} GB/s $\approx$ \val{qflat} ms. (b) Shard \val{qshard} MB;
  \val{qringfactor} $\times$ \val{qshard} MB $\div$ \val{beta} GB/s $\approx$ \val{qinter} ms; NVLink phases
  \val{qnvfrac} of \val{qgb} GB at \val{nvlink} GB/s $\approx$ \val{qrs} ms each; total \val{qtotal} ms,
  \val{qgain}$\times$ faster. (c) Shards take turns: cross-node \val{hiergpus} $\times$ \val{qinter} $\approx$
  \val{qshared} ms, about \val{qsharedtotal} ms in all, near flat (hierarchy card).}
\columnbreak
\ans{3}{It removes \val{qremoved} ms for \val{qcodec} ms: \val{qcompressed} + \val{qcodec} = \val{qnew} ms,
  \val{qspeedup}$\times$ faster. On NVLink it removes only \val{qfastremoved} ms: \val{qfastcompressed} +
  \val{qcodec} = \val{qfastnew} ms, \val{qfastslow}$\times$ slower (payback condition, compression card).}
\ans{4}{Record each rank's GPU, NIC, node and rail; rerun the matched benchmark, changing only
  placement. If it slows, placement; if not, overlap or contention (ring card).}
\end{answers}

\begin{leftcol}
\sheetsec{Worked example}
\begin{practicebox}\boxtitle{Is a \val{bucket} MB gradient bucket bandwidth bound?}
Ring AllReduce of one \val{bucket} MB bucket on the NDR port, on \val{small} and \val{large} ranks.\par\vspace{5pt}
{\centering
\begin{tikzpicture}[sheet]
  \pgfmathsetmacro\s{62}
  \draw[dim] (8,0) -- (232,0);
  \foreach \e/\t in {1/{10 KB},2/{100 KB},3/{1 MB},4/{10 MB}}{
    \draw[dim] ({8+(\e-1)*\s},0) -- ({8+(\e-1)*\s},-2.5);
    \node[small, anchor=north] at ({8+(\e-1)*\s},-4) {\t};}
  \pgfmathsetmacro\xl{8+(\val{logchunklarge}-1)*\s} \pgfmathsetmacro\xn{8+(\val{lognstar}-1)*\s}
  \pgfmathsetmacro\xs{8+(\val{logchunksmall}-1)*\s} \pgfmathsetmacro\xb{8+(\val{logbucket}-1)*\s}
  \draw[accent, dashed, line width=0.7pt] (\xn,-2) -- (\xn,26);
  \node[note, anchor=south] at (\xn,27) {knee, \val{nstar} KB};
  \filldraw[fill=accent, draw=accent] (\xl,0) circle (2.4pt);
  \filldraw[fill=accentsoft, draw=accent, line width=0.5pt] (\xs,0) circle (2.4pt);
  \filldraw[fill=white, draw=muted] (\xb,0) circle (2.4pt);
  \node[note, anchor=south, align=center] at (\xl,5) {\val{large} ranks\\chunk \val{chunklarge} KB};
  \node[note, anchor=south, align=center] at (\xs,5) {\val{small} ranks\\chunk \val{chunksmall} KB};
  \node[note, anchor=south east, align=right] at ({\xb+3},5) {bucket\\\val{bucket} MB};
\end{tikzpicture}\par}\vspace{6pt}
\begin{steps}
\item Knee: $\LatencyBandwidthCrossoverSize$ = \val{alpha}\,$\MicroUnitPrefix$s $\times$ \val{beta} GB/s = \val{nstar} KB.
\item Chunk $\CollectivePayloadSize/\DeviceCount$: \val{chunksmall} KB on \val{small} ranks (above: bandwidth
  bound), \val{chunklarge} KB on \val{large} (below: latency bound).
\item Startup: \val{stepssmall} $\times$ \val{alpha}\,$\MicroUnitPrefix$s = \val{startupsmall} ms;
  \val{stepslarge} $\times$ \val{alpha}\,$\MicroUnitPrefix$s = \val{startuplarge} ms.
\item Bytes about \val{bytessmall} and \val{byteslarge} ms. Ring: \res{\val{ringsmall} ms} and
  \res{\val{ringlarge} ms}, \val{ringslowdown}$\times$ slower for the same bytes.
\end{steps}
\takeaway{The regime belongs to the chunk, not the tensor. On a big group, cut rounds.}
\end{practicebox}

\vspace{10pt}
\sheetsec{Your turn: finish it}
\begin{practicebox}\boxtitle{Would the simple tree beat the ring on \val{large} ranks?}
\textbf{Given:} crossover \val{large} $\times$ \val{alpha}\,$\MicroUnitPrefix$s
$\times$ \val{beta} GB/s $\div$ \val{log2large} = \val{crossover} MB; the tree's \val{treerounds} rounds cost
\val{treestartup} ms of startup, and each sends all \val{bucket} MB in \val{perround} ms.\par
\textbf{You:} find the tree's time, compare it with the ring's \val{ringlarge} ms using the crossover,
then estimate tree rounds with ring bytes.\par\vspace{34pt}
\end{practicebox}
\end{leftcol}%
\begin{rightcol}
\sheetsec{Common mistakes}
\mistake{``An expert exchange costs its average traffic.''}%
  {\val{moegpus} GPUs, \val{moeexperts} experts: each receives
   \val{rxuniform} MB, at least \val{bounduniform}\,$\MicroUnitPrefix$s. One rank at \val{skew}$\times$ the load
   receives \val{rxhot} MB: at least \val{boundhot}\,$\MicroUnitPrefix$s for the whole layer.}{}
\mistake{``Asynchronous AllReduce hides all communication.''}%
  {\val{layers} layers, \val{bwdlayer} ms backward and \val{commlayer} ms reduction each: overlapped
   \val{bwdtotal} + \val{commlayer} + \val{layersless} $\times$ \val{exposed} = \val{overlapped} ms, not
   \val{bwdtotal}; sequential \val{sequential}.}{}
\mistake{``Error feedback is optional for compression.''}%
  {Gradients \val{efgrads}, rounded to whole units: alone, \val{efnaive} of the true \val{eftrue} is sent.
   With error feedback, \val{efsent} is sent and $\val{efresidual}$ is kept.}{}

\vspace{4pt}
\sheetsec{Check yourself}
\begin{qs}
\item Each rank updates only its weight shard, but the code AllReduces the full gradient. What does
  each rank hold, what should it call, and what is saved?
\item \assumed{} \val{qgb} GB on \val{qnodes} nodes of \val{hiergpus} GPUs, NVLink \val{nvlink} GB/s, one
  \val{beta} GB/s port per GPU. (a) Flat ring time? (b) Shard, cross-node phase and hierarchical
  total? (c) And if a node's GPUs share one uplink?
\item \assumed{} a \val{qcomm} ms synchronization; compression gives an effective \val{qratio}$\times$ at
  \val{qcodec} ms of encoding. Does it pay here, and on NVLink where it takes \val{qfast} ms? Which
  condition decides?
\item An AllReduce benchmark reaches full bus bandwidth; the job's reductions get a fraction.
  Fabric or placement? What do you record?
\end{qs}
\end{rightcol}
\end{document}
